\head{第十八版整合核查}
\source{2026年11月25日；基线为第十七版增补第三辑实际源码；新增来源为Dot原批号69--108的四十个实存工程包}

\egasection{一、此次工作与证据边界}
本轮整合的对象为已存译稿，不把四十批已有内容列为本轮从头新增翻译。原第十七版工程、Dot原ZIP和展开原件保持不变；在新的模块工作副本中接入续稿、替换对应回校段，并记录原文件散列、裁剪边界、局部宏适配和迁移去向。

第69--72批接到15.7之后，收入附录17.16的16.1--16.14及相关作者按语；第73--85批为第三篇，含射影丛与多重射影丛、相对除子、伪态射及线性系；第86批为1965年9月29日致迪厄多内书信，保留划删技术段的身份；第87批为法文物理150页提纲和153页异扫核对。第88--108批是第一篇§§1--7、8.1--8.13的回校，不另算新的线性章节。既有8.14及后续正文仍保留。

每批原有的“范围与接续”、独立封面、当时未并入合订本的说明只属于制作记录，未混入新的连续正文。正文中的作者编排意见、删划层和原注仍予保留。长篇译者补证按来源移到下列报告章节，保留完整内容并注明非原文。旧稿特有译注的保留与去向见各组整合记录。

本轮对部分重要订正查看法、英原页图，并检查结构接合。相关原批次的其他证明继续以其原有译校身份收入；并不声称本次重新独立证明全部命题。第I、II、III章、第0章及旧报告没有因合订而获得新一轮全卷审定。

\egasection{二、7.4：严格条件、英译放宽与附加条件版本}
法文物理29页与英文物理20--21页的对应条件必须区别。新增回校记录指出，法文使用严格不等式，英文把它放宽为非严格不等式；旧译所列纯不可分反例针对这个放宽。不能据同一反例断言法文严格条件下的一般命题也已被否定。

另一方面，把英译的不等式改回严格形式，并不自动填补一般固有、平坦情形所需的有限表示与射影切割论证。正文区分有明确附加条件的译校采用版本与法文一般原断言；原断言、原归纳思路、反例和补证在对应报告中保留。这个处理是将不同范围的结论分清，并非为原稿一般版本补发一个无条件的正确性结论。

\egasection{三、光滑二次型：判别式的符号与特征2}
按16.1的极化矩阵及16.5的校正判别式约定，标准二次型
\[
 X_1X_2+\cdots+X_{2m-1}X_{2m}
 \quad\text{或}\quad
 X_1X_2+\cdots+X_{2m-1}X_{2m}+X_{2m+1}^2
\]
的校正判别式为$(-1)^m$。它是单位，但一般不是$1$。第71批在16.9的证明中把原稿所写的$1$改为“单位”，并在脚注交代符号依据。本轮另作了秩1至8的整数矩阵代入核对；这是对所用约定的精确复核，不代替任意秩分类定理的证明。

还核对了三元二次型
\[
 Q=ax^2+by^2+cz^2+dxy+exz+fyz
\]
的极化矩阵：
\[
 \begin{gathered}
 B_Q=\begin{pmatrix}2a&d&e\\d&2b&f\\e&f&2c\end{pmatrix},\\[.4em]
 \frac{\det B_Q}{2}=4abc+def-af^2-be^2-cd^2.
 \end{gathered}
\]
这里先在整数多项式环中作可整除的计算，再将所得多项式用于任意底环；不是在特征2的环内任意除以$2$。对$xy+z^2$，该式为$-1$。而在特征2，极化矩阵行列式为零；这正说明奇秩光滑二次型不能仅用极化型非退化来表述。

秩一型$ax^2$与$by^2$各有校正判别式$a,b$，其直和的偶秩极化行列式为$4ab$。因此奇秩加奇秩的直和公式中的系数$4$不能删去。符号计算的脚本和具体结果保存为工程的\texttt{checks/check\_quadratic\_examples.py}及\texttt{checks/quadratic\_examples.json}。

第16.9--16.14关于忠实平坦、有限表示的局部标准化，与光滑局部标准化应区分；特征2、奇秩时相应轨道映射不必光滑。原批次对此的双数计算作为译校说明保留。本轮的矩阵核对不扩大其原有适用范围。

\egasection{四、相对除子参数与删划草稿补译}
第一篇§1的预告把次数$k$的参数模写成了$\operatorname{Sym}^k(E^\vee)$。在本稿的一维商约定下，方程所在的模是$p_*\OO_P(k)=\operatorname{Sym}^k E$，参数化它的局部直和因子线所需的射影丛为
\[
 \mathbf P\bigl((\operatorname{Sym}^k E)^\vee\bigr).
\]
不能在任意底上无条件交换对偶与对称幂。第十八版在正文订正参数模，脚注保留原式及其依据；这不影响$k=1$的超平面构造，也不把这个预告改写成已在此处证明的一般可表性定理。

本轮另对法文物理126、136、140页的大叉删划块补译可辨文字，作为第三篇之后的单独补遗。原批次仅述及这些删划块而未译其主体；新增补译不计作第69--108批已经存在的翻译量。密集覆盖字、未竟句和难辨边注按原状标明，不恢复为作者有效定稿。

物理136页删稿由$N\simeq\operatorname{pr}_1^*L\otimes\operatorname{pr}_2^*\OO_{P^\vee}(1)$恢复$L$时，原句只写推前$N$。在参数射影丛具有正秩纤维的情形，应当先消去该扭转：
\[
 L\simeq\operatorname{pr}_{1*}
 \bigl(N\otimes\operatorname{pr}_2^*\OO_{P^\vee}(-1)\bigr).
\]
订正及原式写在补译脚注。物理140页划去的旧推论5.3只凭忠实平坦、局部有限表示就断言参数丛到除子函子为单态射；第三篇4.3已给出这种弱假设不足的反例。因此删稿按其被划去且不能直接采用的身份译存，较强条件来自4.3(a)，不冒称旧推论原本已经具备。

\egasection{五、底本末页与真实接续}
已实际打开两份完整底本，英文编辑稿为112页，末页为致迪厄多内书信及编者脚注82--84；法文扫描为171页，末页仍是光滑二次型附录的末段，论述有限局部自由基变换与结构群的约化。两份载体的末页内容不同，不能以英文目录顺序猜测法文扫描末页。

第十八版将这批现存后续材料接入后，英文物理78--112页不再是独立于合订本的待并入部分。法文物理150页手写提纲及153页异扫仍保留难辨字的限定；物理153页与卷首同文不在正文重复排印。到达载体末页不能证明全部法文页均已逐行校勘。

第一篇最新回校点为8.13（法文物理45--46页，英文31--32页）。下一回校从法文46页下部“Combinant 8.13”及8.14开始，接47页；已有8.14及后文作为对照基线继续保留。正式第IV章仍是独立未完成分支，从\texttt{EGAIV.1.pdf}物理221页提要、222--223页章首说明及223页（1.1.1）接续。

\egasection{六、未决状态与技术核验的用途}
8.7(c)缺续、部分覆盖字与边注、9.4拟写提示和9.5残句等不能靠编号完整或编译成功补齐。第8.13的回校采用状态也不直接审定8.14的充要方向。一般固有情形的7.4、既有13.3--13.4和14.4(a)等问题须按各条实际条件继续处理。具体更新列于工程\texttt{PENDING\_ISSUES.md}，旧状态在原归档中保留。

技术检查分别核对原件保存、实际TeX输入、重叠段替换、单文件展开、四份实际XeLaTeX日志、模块与单文件PDF逐页文字／尺寸／像素一致性，以及变动页的渲染抽查。最终运行结果和实际查看的页码以工程\texttt{checks/}中的本版记录为准。逐页自动比对只比较两种构建的输出，不能说明每页数学文字已经人工审完；抽样渲染只覆盖明确列出的页面。
